% !PRINTED 572
\weberpage{572}

%<<CONT
\weq{1}{
\begin{aligned}
&\text{a)} \quad \tau(u) = \frac{g_2 g_3}{G} \wp(u), \quad \text{in general,} \\
&\text{b)} \quad \phantom{\tau(u)} = \frac{\wp(u)^2}{g_2}, \quad \text{if } g_3 = 0, \\
&\text{c)} \quad \phantom{\tau(u)} = \frac{\wp(u)^3}{g_3}, \quad \text{if } g_2 = 0.
\end{aligned}
}

And if then $\mu$ is, as above, an integer of the field $\Omega$, we can set in all three cases:

\weq{2}{\tau(\mu u) = \frac{\Psi(\tau)}{\Phi(\tau)},}

in which $\Phi, \Psi$ are integral functions of $\tau$ which, apart from constant factors, agree in case (1) a) with the $R, P$ of the preceding section, and in the cases (1) b) and (1) c) arise from these by raising to the square or to the cube [§\,152, 4), 5)].

We arrange the fraction (2) so that the highest power of $\tau$ in the denominator $\Phi$ has the coefficient 1. The remaining coefficients of $\Phi$ are then rationally composed of quantities of the form

\[
\tau\left(\frac{\nu}{\mu}\right) = \tau\left(\frac{\nu \mu'}{m}\right),
\]

and since $\nu\mu'$ is a period of $\tau$, these quantities belong to the roots of the division equation of the periods, and are therefore, since the modulus of the corresponding elliptic functions is an algebraic number, themselves algebraic numbers (§\,58, §\,61).

From this it follows:

\begin{quote}
2. \gesperrt{The coefficients in $\Phi(\tau)$ are algebraic numbers.}
\end{quote}

We extend the fraction $\Psi : \Phi$ by multiplication of numerator and denominator by the product of the conjugate values $\Phi', \Phi'', \Phi''', \ldots$ and obtain:

\weq{3}{\tau(\mu u) = \frac{Z}{N},}

in which now $Z$ and $N$ are integral functions of $\tau$ which may contain a common divisor, while $N$ has rational numerical coefficients.
